\section{Introduction}

Spin correlations of unstable particles are measured through the angular distributions of their decay products. On this basis entanglement has been reported for top-quark pairs at the LHC \cite{atlas2024,cms2024,cms2024ljets}, following the proposal of Ref.~\cite{afik2021}, and a wider programme has formed around top quarks, $\tau$ leptons, hyperons and massive gauge bosons \cite{barr2024,ashby2023,tornqvist1981,besiii2025}.

It has been known since the first proposals that such data cannot test local realism \cite{abel1992}. Each particle is measured by its own decay, that is, with a single setting; all recorded quantities commute; and the recorded momenta themselves constitute a local hidden-variable model. The point has recently been restated and sharpened \cite{abel2025,bechtle2025,low2025}, and what can still be inferred has been discussed in Refs.~\cite{li2024,fabbrichesi2025bell}. What the measured coefficients do establish is entanglement of the reconstructed spin state, given the quantum description of the decay and in particular the value of its spin analysing power.

This paper addresses a question that is well posed for one fixed measurement. A decay is a rotation-covariant positive-operator-valued measure (POVM) on the spin of the parent, with a continuum of outcomes. Many functions of the outcome define the same effect: the angular distribution of a spin-$j$ parent contains spherical harmonics of degree at most $2j$, so every function of the recorded direction that is orthogonal to these integrates to the zero operator. Such operational equivalences are what a generalized-noncontextual ontological model \cite{spekkens2005} must respect: equivalent preparations are represented by the same distribution over ontic states, equivalent effects by the same response function. The existence of such a model is a notion of classical explainability with several equivalent forms: simplex embeddability of the accessible fragment \cite{schmid2021,gitton2022,selby2023frag}, which is decidable by a linear programme \cite{selby2024lp}, and nonnegativity of some quasiprobability representation \cite{spekkens2008}. It can fail for a single measurement \cite{selby2023incomp}, and for a bipartite state with all local measurements available it fails exactly when the state is entangled \cite{plavala2024}. The recorded-outcome model of Refs.~\cite{abel1992,abel2025} is not of this kind: it assigns different response functions to equivalent effects (Sec.~\ref{sec:models}).

We determine exactly when decay data admit a noncontextual model.

(i) \emph{Spin $1/2$, one particle.} With polarisation $P$ and analysing power $\kappa$ a model exists if and only if $P|\kappa|\le 1/3$ (Theorem~\ref{thm:half}).

(ii) \emph{Two spin-$1/2$ particles.} Write the joint distribution of the two recorded directions as $(4\pi)^2p=1+\ba\!\cdot\!\bq_A+\bb\!\cdot\!\bq_B+\bq_A^{\T}c\,\bq_B$. A model exists if and only if the operator $\rhoraw=\tfrac14(\openone+\ba\!\cdot\!\bsig\otimes\openone+\openone\otimes\bb\!\cdot\!\bsig+\sum_{ik}c_{ik}\sigma_i\otimes\sigma_k)$ is positive semidefinite and has a positive partial transpose (Theorem~\ref{thm:pair}). The test uses the raw coefficients; the entanglement test uses the coefficients divided by $\kappa_A\kappa_B$. Contextual data imply an entangled state, and the converse holds for all states only if $|\kappa_A\kappa_B|=1$.

(iii) \emph{Spin $1$.} A decay is specified by helicity weights $(w_+,w_0,w_-)$, equivalently by the vector and tensor ratios $\eta_1=w_+-w_-$ and $\eta_2=1-3w_0$. On all spin-$1$ states a model exists if and only if $(\eta_1,\eta_2)$ lies in the convex hull $\cR$ of the ellipse $3\eta_1^2+25(\eta_2-\tfrac15)^2\le1$ and the point $(0,-\tfrac45)$ (Theorem~\ref{thm:one}). For a spin-singlet pair with one decay on each side the products of the ratios must lie in a second region $\cR_2$ (Theorem~\ref{thm:onepair}).

(iv) \emph{Spin $j$.} The largest vector analysing power of a noncontextual polarimeter is the largest node of the Gauss--Legendre rule with $j+1$ nodes (integer $j$) or of the Gauss--Radau rule with $j+\tfrac32$ nodes (half-integer $j$): $\xi_{2j}=1/3$, $1/\sqrt3$, $(1+\sqrt6)/5$, $\sqrt{3/5}$ for $j=\tfrac12,1,\tfrac32,2$ (Theorem~\ref{thm:spinj}). For two particles of arbitrary spins a noncontextual model is a local model whose response functions are nonnegative polynomials of degree $2j$ in the recorded direction (Proposition~\ref{prop:pairpoly}), which implies $\|\langle\bq_A\bq_B^{\T}\rangle\|_1\le\xi_{2j_A}\xi_{2j_B}$ for the matrix of first moments, whatever the state and the decays.

(v) \emph{Data.} Section~\ref{sec:apps} applies these statements to the decays of muons, $\tau$ leptons and top quarks, to the top-pair coefficient $D$ measured by ATLAS and CMS, to the BESIII parameters of $J/\psi\to\Lambda\bar\Lambda$, and to decays of vector bosons and vector mesons.

\emph{Relation to earlier work.} Fucci \emph{et al.}~\cite{fucci2026} construct Bell tests for decay processes in which independently sampled analysis axes play the role of settings and detector regions give binary or null outcomes; a violation requires postselection, and they show that a local model exploiting it must be measurement contextual, with an inequality violated when the effective analysing powers are large enough. Here no auxiliary settings and no postselection are introduced: the object is the fragment formed by the recorded directions themselves, and the conditions are necessary and sufficient. Fabbrichesi \emph{et al.}~\cite{fabbrichesi2025pra,fabbrichesi2025prd} evaluate Kochen--Specker-type inequalities on spin density matrices reconstructed from collider data; the measurements entering those inequalities are projective spin measurements evaluated on the reconstructed state, not the decay. Li, Shen and Yang~\cite{li2024} exclude a class of hidden-variable models with classical spins and known analysing powers. The equivalence of entanglement and contextuality for all local measurements \cite{plavala2024} becomes a strict implication here because a decay supplies one measurement. For a single qubit our threshold is the entanglement-breaking threshold of the depolarising map \cite{horodeckiEB2003,ruskai2003}; noise thresholds for preparation noncontextuality alone are weaker \cite{marvian2020}, and the classification of individual processes by generalized noncontextuality is developed in Ref.~\cite{zhang2025prx}. Related notions of operational contextuality, including contextuality revealed by a single measurement, are studied in Refs.~\cite{jokinen2024,jokinen2026}.

\emph{The graph.} The question arose from two released concept graphs, each with a state, a model and an observable layer. Their convention is that a relation known in the literature is an edge; the charged-current vertex of the particle-physics graph \cite{genhep} and the generalized-contextuality node of the Nonclassicality Network \cite{gengpt} were not connected. The statements of this paper are the routes that connect them, through new nodes for the decay polarimeter, the steered ensembles and the tests. Each route has a grade (exact, logic, numeric, import, definition, or written proof) and, where possible, an executable certificate; Sec.~\ref{sec:graph} and Appendix~\ref{app:catalogue} give the encoding and the results of the checks.
